% TOP_PROBLEM: {"id":30007750,"problem_number":"LOCAL-30007750","title":"Remote work and city economies","category_id":21,"category":{"id":21,"name":"economics","display_name":"Economics","description":"Open economic theory statements, published research agendas, and proposed scoped specifications; question provenance and historical priority are qualified per record.","slug":"economics","order_index":21,"created_at":"2026-10-08T00:00:00Z"},"status":"provenance_pending","external_url":null,"legacy_ids":[],"published":false,"statement":"Estimate the twenty-year metropolitan welfare effect of a specified permanent increase in home-working days after firms, residents, rents, and public services adjust.","economics":{"original_id":239,"field":"Urban, housing and spatial economics","kind":"empirical","formalization_basis":"proposed_specification","source_status":"not_verified","priority_status":"not_established","priority_note":"No readable primary passage explicitly stating this precise open question was verified. A supporting paper, abstract, or topic citation is insufficient to establish origin or unresolved status.","earliest_verified_reference":null,"current_reference":null,"formalization_note":"Proposed scoped research specification. The original catalogue prompt was a synthesis, and no canonical conjecture or verified original formulation is claimed. The supporting source, when retained, establishes context or existing results rather than this exact open question. Mathematical scope was checked independently; added definedness, feasibility, or identification conditions belong to this proposed formalization."},"statusQualification":"Provenance pending: an explicit published statement of this exact open question was not verified. This is a catalogue-authored candidate specification, not a certified open conjecture. Proposed scoped research specification. The original catalogue prompt was a synthesis, and no canonical conjecture or verified original formulation is claimed. The supporting source, when retained, establishes context or existing results rather than this exact open question. Mathematical scope was checked independently; added definedness, feasibility, or identification conditions belong to this proposed formalization.","statusReviewedAt":"2026-10-08"}
% FIRST_POSTED: 2026-10-08
% ENTRY_KIND: statement_only
% !TeX program = lualatex
\documentclass[11pt]{article}
\usepackage{fontspec}
\usepackage[a4paper,margin=25mm]{geometry}
\usepackage{amsmath,amssymb,mathtools}
\usepackage{xurl}
\usepackage[hidelinks]{hyperref}
\newcommand{\sref}[2]{\hyperlink{source:#1:#2}{[S#2]}}
\setlength{\emergencystretch}{3em}
\begin{document}
% problemId:30007750
\section*{Remote work and city economies}\label{30007750}
\subsection{Definitions and mathematical statement}
Consider a metropolitan economy with residential neighborhoods, an office center, and occupations differing in remote-work feasibility. Let \(a=1\) denote a permanent specified increase in permitted home-working days and \(a=0\) the baseline commuting regime. Firms and workers choose location, office demand, employment, and commuting; land rents and local tax bases adjust. For resident \(i\), define \(v_i(a)\) as twenty-year monetary-equivalent welfare. Let \(\Pi_f(a)\) be discounted firm profit and \(K(a)\) public infrastructure and service resource cost. Define
\[\Delta W=E\!\left[\sum_i\{v_i(1)-v_i(0)\}+\sum_f\{\Pi_f(1)-\Pi_f(0)\}\right]-\{K(1)-K(0)\}.\]
Household monetary equivalents exclude profit distributions already entered in the separate profit term. Estimate how agglomeration productivity, commuting savings, office conversion, and municipal service responses contribute to this contrast. Specify the population accounting boundary, including residents who leave, and treatment of property transfers. Policy-driven variation in remote-work permission must be separated from occupation demand shocks and worker selection. Short-run office-price movements do not identify the full long-run welfare effect. The original review discusses remote work and urban real estate, but an explicit open passage for this exact twenty-year equilibrium estimand was not verified. This remains a proposed empirical model comparison with no claimed original priority.

\par\textbf{Formulation and scope.} Proposed scoped research specification. The original catalogue prompt was a synthesis, and no canonical conjecture or verified original formulation is claimed. The supporting source, when retained, establishes context or existing results rather than this exact open question. Mathematical scope was checked independently; added definedness, feasibility, or identification conditions belong to this proposed formalization.

\par\textbf{Open-question provenance.} An explicit published open-question statement was not verified. Sources below are supporting literature and must not be treated as a first listing.

\par\textbf{Historical priority.} No readable primary passage explicitly stating this precise open question was verified. A supporting paper, abstract, or topic citation is insufficient to establish origin or unresolved status.

\subsection{Short English statement}
Estimate the twenty-year metropolitan welfare effect of a specified permanent increase in home-working days after firms, residents, rents, and public services adjust.

\subsection{Sources}
\begin{itemize}\raggedright
\item[\textnormal{[S1]}]\hypertarget{source:30007750:1}{} Year not verified. The Remote Work Revolution: Impact on Real Estate Values and the Urban Environment (2022). Suggested supporting reference from the initial catalogue; exact open-question passage not verified..
\url{https://www.nber.org/papers/w30662}.
{\small\textit{Unverified bibliographic lead Bibliographic lead only. It is not evidence of an original explicit open-question statement.}}
\end{itemize}
\end{document}
